\begin{workedexample}[Worked Example: Colpitts frequency]
A Colpitts oscillator uses a tank inductor $L = 10\ \text{nH}$ with a
capacitive divider $C_1 = C_2 = 2\ \text{pF}$. The feedback capacitors appear in
series, $C_s = C_1C_2/(C_1+C_2) = 1\ \text{pF}$, so
\[ f_0 = \frac{1}{2\pi\sqrt{L\,C_s}} = \frac{1}{2\pi\sqrt{(10\times10^{-9})(1\times10^{-12})}} = 1.59\ \text{GHz}. \]
The divider ratio sets the feedback needed to sustain oscillation.
\end{workedexample}

\begin{keytakeaways}
\begin{itemize}[leftmargin=1.4em,itemsep=2pt]
\item Start-up: $|R_d| > R_{res}$; steady state: $|R_d| = R_{res}$ (nonlinear limiting).
\item Colpitts: $f_0 = 1/(2\pi\sqrt{L C_s})$ with $C_s = C_1C_2/(C_1+C_2)$.
\item Crystals give huge $Q$ --- best stability and lowest phase noise.
\end{itemize}
\end{keytakeaways}

\begin{exercises}
\item A Colpitts tank has $L = 5\ \text{nH}$, $C_1 = C_2 = 4\ \text{pF}$. Find $f_0$.
\item Why buffer an oscillator's output?
\item Why does a crystal oscillator have far lower phase noise than an LC one?
\end{exercises}

\furtherreading{Lee~\cite{lee} (ch.~14--15); Razavi~\cite{razavi} (ch.~8).}
